\smalldoc{The Coleman-Mandula theorem}{CM thm}


The Coleman-Mandula theorem 
concludes that the only possible Lie algebra that generates symmetries of the $S$ matrix is locally isomorphic
to a direct sum, where one of the summands is the Poincaré algebra. The theorem basically shows that
spacetime symmetries do not mix with “internal” ones. (However, if we relax
the condition that the symmetry algebra be an ordinary Lie algebra, this allows for more possibilities: namely, supersymmetry.)

 Whatever the quantum field theory at hand is, it has a Hilbert space that looks like $\mathcal{H} := \mathcal{H}^{(1)} \oplus \mathcal{H}^{(2)}\oplus \dots$, where $\mathcal{H}^{(n)}$ is the Hilbert subspace describing the possible states occupied by $n$-particles, with  $\mathcal{H}^{(n)}\simeq  \bigotimes^n_{i=1} \mathcal{H}^{(1)}$.
   The \textit{$S$ matrix} is a unitary (probability-preserving) operator $S: \mathcal{H}\to\mathcal{H}$ which acts on asymptotically-free input states (called \textit{in-states}) and gives asymptotically-free output states (called \textit{out-states}). In other words, if you feed the $S$ matrix some particles in a scattering experiment, it gives you the probability of what happens after the interaction (scattering) takes place.  

A \textit{symmetry group} $G$ of $S$ is a group such that its infinitesimal transformations (generators) are unitary operators $U:\mathcal{H}\to\mathcal{H}$ which do not mix the fixed-particle subspaces and which commute  with $S$. If the Poincar\'e group $P$ is a symmetry group of $S$, then we say    the theory is \textit{Lorentz-invariant}. In this case the asymptotic one-particle states are classified by the irreducible representations of $P$, so we label the states like   $\ket{\alpha,\lambda,p}$, where $p$ is on-shell momentum, $\lambda$ is spin/helicity, and $\alpha$ is any additional discrete particle labels.  An \textit{internal symmetry} of $S$ is one that commutes with $P$.  Of course, we can extract the scattering part of the $S$ matrix by subtracting off the identity; we write $S = 1 - i(2\pi)^4 \delta(P_\mu - P_\mu') T$, where the elements of $T$ are the scattering amplitudes. Now we are ready to state the theorem: 

\begin{theorem}
     Let $G$ be a connected symmetry group for the $S$-matrix and assume the following five conditions: \begin{enumerate}[\bfseries 1.]
    \item (Lorentz invariance)  $G$ contains a subgroup locally isomorphic to $  P$. 
    \item (Particle-finiteness)  All particle types correspond to positive-energy reps of $P$. For any $\varepsilon>0$, there are only finitely-many particles with mass $0 < m< \varepsilon$. 
    \item (Weak elastic analyticity)  Elastic scattering amplitudes are analytic functions of center-of-mass energy $s$ and momentum transfer $t$ in some neighbourhood of the physical region, except at normal thresholds (in other words, almost everywhere). 
    \item (Occurrence of scattering)  Let $\ket{p}$ and $\ket{p'}$ be 1-particle momentum states and let $\ket{p,p'}$ be built from these. Then $T\ket{p,p'}\neq 0$ almost everywhere.
    \item (Ugly technical assumption)   $G$ has well-defined infinitesimal generators on 1-particle states, in the sense that for every  one-parameter subgroup $g(t)\subset G$ with $g(0) = e$, the derivative  $$ \frac{d}{dt}\Bigg|_{t=0}\braket{x,g(t) y}    $$ exists and is expressable as  $ \frac{d}{dt}\Big|_{t=0}\braket{x,g(t) y}  := i\braket{x,Ay}$ for some distribution $A$. (For example, $g(t) = e^{it A}$, but this is stronger than necessary).
\end{enumerate} \textit{Then}, $G$ is locally isomorphic to $P$ times an internal symmetry group, or equivalently  $\mathfrak{g}\simeq\mathfrak{p}\oplus \mathfrak{internal}$ locally (where $\mathfrak{internal}$ denotes the Lie algebra of the internal symmetry group).  
 \end{theorem} 

 \begin{proof}
     We will  show that $A$ is the sum   of an infinitesimal
Lorentz transformation, an infinitesimal translation,
and an infinitesimal internal symmetry transformation.  For simplicity denote $A(p',p) := \braket{\alpha',\lambda',p'|A|\alpha,\lambda,p}$  (suppressing the spin and particle type indices). For two-particle states, $A$ readily extends to \begin{align}
    A(p',p,q',q) = A(p',p) \otimes \delta(q' - q)\mathbbm{1} +  \delta(p' - p)\mathbbm{1}\otimes A(q',q) 
\end{align} where $\mathbbm{1}$ denotes the identity with respect to the suppressed indices. We do the proof in multiple steps. 
\vspace{3pt}


\noindent\textbf{Step 1.} \textit{Show that $A(p',p)$ vanishes everywhere except when $p' = p$; in other words, $A$ has no matrix elements between different 4-momenta.}
\vspace{2pt}


To this end, we define a momentum projection $A_f:= \int d^4a e^{i p \cdot a}Ae^{-ip\cdot a} \tilde{f}(a)$, where $f$ is an arbitrary test function with support in some region $R$, and $\tilde{f}$ is its Fourier transform. The distribution $A_f$ is still a generator of $\mathfrak{g}$ since it is just modified by weighted translations. Computing the matrix elements, one easily obtains that $A_f(p',p) = f(p' - p)A(p',p)$, which means that $A_f$ is only supported in the region where the difference $p' - p$ remains inside $R$.   This is significant: because of the particle-finiteness assumption, the one-particle state momenta are only supported on a set of countable mass hyperboloids; i.e. $p^2 = m^2_i$, with at most $i \in \mathbb{Z}$. This means that, for fixed $r\in R$, the system of two equations  $p^2 = m_i^2$  and $(p + r)^2 = m_j^2$ describes a slice of a mass hyperboloid, and varying $r$ by a small amount turns this slice into a thin band $L$. By definition, everything outside the band but within the hyperboloid is annihilated by $A_f$. So if $R$ is small enough, then for any $p\in L$ and any $r\in R$, $p + r$ takes us outside the band and hence the corresponding state is annihilated by $A_f$. 

First we will consider the lowest mass hyperboloid. Let $p \in L$ and $q,p',q'\notin L$ but still within the hyperboloid, such that $p + q  = p' + q'$. Further suppose that $s = (p + q)^2$ is small enough so that we are not taken to the next hyperboloid (in other words, the energy involved is not large enough to produce a   pair of particles of the next type). Since $A_f$ is a generator of $\mathfrak{g}$, it commutes with $S$, so we have \begin{align}
    \bra{p',q'} S\Big(A_f\ket{p}\otimes \ket{q} + \ket{p}\otimes A_f\ket{q}\Big) = \Big( \bra{p'}A_f\otimes \bra{q'} +  \bra{p'}\otimes \bra{q'}A_f\Big)S\ket{p,q},
\end{align} but only $A_f\ket{p}\neq 0$ so we obtain $ \bra{p',q'} S\big(A_f\ket{p}\otimes \ket{q} \big) = 0$. But then this can only happen if $S =0$, which is possible for some values of $s$ and $t = (p - p')^2$. However,  because we can move $q,p',q'$ around by rotating them, we actually obtain that $S = 0$ for a continuous range of $s$ and $t$, which by analyticity implies $S \equiv 0$ for scattering on this hyperboloid. This contradicts our assumption and thus $A_f = 0$ for this hyperboloid. Applying the same argument to the rest of the hyperboloids gives that $A_f = 0$ everywhere, and since $f$ was arbitrary except that it had support away from zero, we conclude that $A$ must have support nowhere except at $ p = p'$, as desired.

\vspace{3pt}
\noindent \textbf{Step 2.} \textit{Rewrite $A(p',p)$ as a distribution in the  momentum variables.} 
\vspace{2pt} 



The technical assumption means that $A$ is a distribution, and from the previous step we can therefore write it as a finite sum of derivatives of  $\delta(p' - p)$: \begin{align}
    A^{(N)}(p',p)=  \sum_{n = 0}^N A^{(n)}_{\mu_1\dots\mu_n}(p) \frac{\partial}{\partial p_{\mu_1}}\dots \frac{\partial}{\partial p_{\mu_n}} \delta( p' - p).
\end{align} We also have (using simple delta manipulations) \begin{align}\label{AAA}
     A^{(N)}_{\mu_1\dots\mu_n}(p) \delta( p' - p) = \frac{1}{N!}\braket{p'|\Big[P_{\mu_1},\big[P_{\mu_2},\dots [P_{\mu_n},A^{(N)}]\big]\Big]|p},
\end{align}  which contains no derivatives of delta functions, hence it is a generator for $N = 0$ representations. We also have the condition  $    [A^{(N)}, P^\mu P_\mu]  = 0 $ which enforces that the momenta are on-shell (the derivatives can only be taken tangent to the hyperboloid $p^2 = m^2$). 

\vspace{3pt}
\noindent\textbf{Step 3.} \textit{Isolate the $N = 0$ representations.}
\vspace{2pt}

Look at the case $N = 0$, i.e. $A^{(N)}(p',p)  = A^{(0)}(p)\delta(p'- p)$. Also writing $ \braket{p',q'|S |p,q} = \delta( p + q - p' - q')S(p,q,p',q') $ (still suppressing non-momentum indices) , then the relation $\braket{p',q'|SA^{(0)}|p,q} = \braket{p',q'|A^{(0)}S|p,q}$ becomes \begin{align}
    S(p,q,p',q') \underbrace{\Big(A^{(0)}(p)\otimes I(q) + I(p)\otimes A^{(0)} (q)\Big)}_{A_{\text{in}}} =  \underbrace{\Big(A^{(0)}(p')\otimes I(q') + I(p')\otimes A^{(0)} (q')\Big)}_{A_{\text{out}}}S(p,q,p',q') ,
\end{align} where $I$ denotes the identity with respect to non-momentum indices at the given momentum. Since $S$ is unitary, then $A_{\text{out}} = SA_{\text{in}}S^{-1} $ is a similarity transformation, giving that $A_{\text{out}}$ and $A_{\text{in}}$ have the same trace. If we denote $N(m_p) := \operatorname{tr}I(p)$, which is the number of particles in the state with momentum $p$ (finite by assumption), then \begin{align}
    \frac{\operatorname{tr} A^{(0)}(p')}{N(m_{p})} +  \frac{\operatorname{tr} A^{(0)}(q')}{N(m_{q})} =  \frac{\operatorname{tr} A^{(0)}(p)}{N(m_{p})}  +  \frac{\operatorname{tr} A^{(0)}(q)}{N(m_{q})} .
\end{align} The only constraint we have is $p + q = p' + q'$, so the above is solved only with $\operatorname{tr} A^{(0)}(p) = (a_\mu p^\mu + b) N(m_p)$, where $a_\mu$ and $b$ are independent of momenta, spin and particle type indices. It follows that  \begin{align}
    A^{(0)} = a_\mu P^\mu + b I + B
\end{align} where $B$ is a traceless $N = 0$ generator. 


\vspace{3pt}
\noindent\textbf{Step 4.} \textit{Show that such $B$ only generates internal symmetries.}
\vspace{2pt}

First observe that such $B$'s yield traceless commutators, so they for their own subalgebra. Pick a basis of them $B_\alpha$ and write $[B_\alpha,B_\beta ] = i C^{\gamma}_{\alpha\beta}B_\gamma$ for some (real) structure constants $C^{\gamma}_{\alpha\beta}$. Let $\Lambda$ be a Lorentz transformation and $U(\Lambda)$ be the corresponding unitary operator. Since conjugation by $U(\Lambda)$ commutes with momentum and preserves the symmetry algebra, we have that $U(\Lambda) B_\alpha U^\dagger(\Lambda)$ is still in the $B$ subalgebra, so we can write $U(\Lambda) B_\alpha U^\dagger(\Lambda) = D^{\alpha'}_{\alpha }(\Lambda) B_{\alpha'}$ for some matrices $ D^{\alpha'}_{\alpha }(\Lambda)$. Since they preserve the commutator we have $C^\gamma_{\alpha\beta} =  D^{\alpha'}_{\alpha}(\Lambda) D^{\beta'}_{\beta}(\Lambda) D^{\gamma}_{\gamma'}(\Lambda^{-1})C^{\gamma'}_{\alpha'\beta'}$, and a similar relation follows for the Lie algebra metric $g_{\alpha\beta}:= C^\gamma_{\alpha\delta}C^\delta_{\gamma\beta}$,  \begin{align}\label{gggg}g_{\alpha\beta} =  D^{\alpha'}_{\alpha}(\Lambda) D^{\beta'}_{\beta}(\Lambda)  g_{\alpha'\beta'}.\end{align} If the subalgebra were to contain momentum generators, since $[B_\alpha,P^\mu] = 0$ we have that the structure constants with momentum indices must vanish, and the metric becomes block-diagonal with only a non-vanishing block $g_{\alpha\beta}$ where the indices are non-momentum. This matrix $g$ is positive-definite and symmetric, so we can write the above condition \eqref{gggg} as \begin{align}
    \left(g^{-1/2}D g\right)  \left(g^{-1/2}D g\right)^\top = 1.
\end{align} This shows that, after a change of basis, $D(\Lambda)$ is orthogonal (unitary), but the only finite-dimensional unitary representations of the Lorentz group are trivial, hence $D(\Lambda) = 1$. This implies that the subalgebra of $B$'s cannot contain momentum indices, and moreover it cannot contain spin indices since the subalgebra commutes with rotations and hence cannot effectuate changes of spin. Thus the $B$'s can only change particle type (``internal'') indices. 



\vspace{3pt}
\noindent\textbf{Step 5.} \textit{Show that there are no generators with $N > 1$.}
\vspace{2pt}

Recall \eqref{AAA}, which is an $N =0$ generator. From Step 3 we can write its trace excluding the delta function as \begin{align}
    \operatorname{tr} A^{(N)}_{\mu_1\dots\mu_n}(p) = a_{\nu\mu_1\dots \mu_N} p^\nu + b_{\mu_1\dots\mu_N},
\end{align} where $a_{\nu\mu_1\dots \mu_N}$ and $b_{\mu_1\dots\mu_N}$ are symmetric in the $\mu_i$ indices. Now consider $ \big[P_{\mu_2},\dots [P_{\mu_N},A^{(N)}]\big] $, which is still an element of the symmetry algebra, and commute it with $P^\lambda P_\lambda$ and take the trace: it must vanish, so we have  \begin{align}
    0 &= \frac{1}{N!}\operatorname{tr} \braket{p'|\Big[P^\lambda P_{\lambda},\big[P_{\mu_2},\dots [P_{\mu_N},A^{(N)}]\big]\Big]|p} = 2p^\lambda \big(a_{\nu\lambda\mu_2\dots \mu_N} p^\nu + b_{\lambda\mu_2\dots\mu_N}\big) \delta(p' - p).
\end{align} This must hold for all $p$ so we get that $b_{\lambda\mu_2\dots\mu_N} = 0$ and $a_{\nu\lambda\mu_2\dots \mu_N}$ is antisymmetric in its first two indices, which by symmetry of $a$ only leaves $a = 0$. This argument relies on there being at least three indices for $a$, hence this holds for all $N > 1$ generators. 


\vspace{3pt}
\noindent\textbf{Step 6.} \textit{Identify the Poincar\'e  and internal parts of the  $N = 1$ generators.}
\vspace{2pt}


For the case $N = 1$ we have \begin{align}\label{refrjeiahf}
       A^{(1)} (p,p') = \Big((a_{\nu\mu } p^\nu + b_{\mu }) \frac{\partial}{\partial p_\mu} +  (\text{$N = 0$ part})\Big)\delta(p' - p),
\end{align} and from the previous step we conclude  $a_{\nu\mu}$ is antisymmetric and $b_\mu $ vanishes. Then the term containing the derivative is just the momentum part of a Lorentz generator. What about the spin part? If we subtract off $a_{\nu\mu}M^{\nu\mu}$ from \eqref{refrjeiahf}, the derivative part cancels and we obtain an $N = 0$ generator. From steps 3 and 4 we conclude that this result must be the sum of a translation and an internal symmetry. 

\textit{Thus}, any $A$ is the sum of a Poincar\'e generator and a generator of internal symmetry, completing the proof.
 \end{proof}
